Quantum error correction (QEC) decoders are frequently described as a real-time systems problem: a decoding decision must be available within the coherence budget of the ancilla qubits, commonly cited as under one microsecond for superconducting devices. We revisit this claim for the smallest interesting case, distance-3 stabilizer codes, on a commodity FPGA accelerator card, and find that the decoder itself is not the bottleneck. We implement three QEC decoder kernels in Vitis HLS for the Xilinx Alveo U55C: a 3-bit repetition code, the Shor \([[9,1,3]]\) code, and a from-source reconstruction of a Steane \([[7,1,3]]\) CSS decoder supporting a lookup-table (LUT) mode and two small-code-specialised decision circuits inspired by minimum-weight perfect matching and union-find (neither is a general graph decoder at this distance; both collapse to the same syndrome-to-qubit bijection the LUT already exploits). Resynthesising the original Shor kernel reproduces its claimed critical path (2.318\,ns) and resource footprint (1~BRAM\textsubscript{18K}, 190 flip-flops, 228 LUT6) exactly, but its claimed five-cycle pipeline latency does not reproduce: the same source, tool, part, and clock target instead report a one-cycle pipeline. We identify and quantify the actual cost of a working hardware read-back path: adding the one-element AXI master output buffer needed to read a result back without BAR4-level host privileges increases pipeline depth by 35 to 70\(\times\) (Shor: 1 to 70 cycles; Steane: 2 to 71 cycles) while the combinational decode logic itself remains under 0.1\% of the device's LUT fabric. \textbf{We built this fixed kernel end to end, placed and routed it, and verified it on a live Alveo U55C}: all 27 single-qubit Pauli self-test cases pass with the result read back through a genuine AXI buffer object, no software fallback, reproduced on two independent runs, and matching the originally claimed syndrome and correction values exactly. The same hardware run gives, for the first time in this line of work, a real post-route resource report (0~BRAM\textsubscript{18K}, 1245 flip-flops, 1018 LUT6, explaining a prior BRAM-count discrepancy as LUTRAM inference) and a real post-route timing closure report (300\,MHz achieved, 0.003\,ns worst negative slack, zero failing endpoints). An opportunistic, non-rigorous latency measurement over 10,000 shots gives a Python-driven host-to-FPGA round trip of 49.5\,\textmu{}s median, 61.0\,\textmu{}s at the 99th percentile. A from-scratch software mirror of each kernel, cross-checked bit-for-bit against an independently written reference decoder over the full input space, gives real Monte Carlo logical-error-rate data at \(10^7\) shots per point: a repetition-code curve that matches its analytic prediction closely, and a Shor code logical error rate under IID-depolarising noise that is an order of magnitude higher than previously reported (15.1\% versus a prior 1.5\% claim at physical error rate 0.10, a gap far outside the statistical uncertainty of either measurement). A two-proportion hypothesis test across the Steane kernel's three decoder modes supports the claim that they are statistically indistinguishable at this code distance. \textbf{The hardware result above covers the Shor kernel only}: the repetition-code and Steane kernels have not yet been placed and routed, and the latency figure is a preliminary Python-loop measurement, not the rigorous, tail-resolving measurement a real-time claim would require. We argue that the useful, defensible claim for a distance-3 FPGA decoder is not "real-time QEC" but a measured characterisation of where the real cost sits: the decode logic is close to free and, now confirmed on real hardware, correct; the host-interface and read-back path is where a genuine real-time budget is actually spent.
